\section{Выводы}
Выполнив вторую лабораторную работу по курсу \enquote{Информационный поиск}, я
написал поискового робота --- компонент обкачки документов. Основное, что я
узнал нового:

\begin{itemize}
    \item устройство краулера на практике: seed-список URL, нормализация URL,
    вежливые задержки между запросами, обработка robots.txt;
    \item хранение документов в MongoDB: схема документа с нормализованным
    URL, сырым текстом, источником и датой обкачки, индексы по url и хэшу
    содержимого;
    \item возобновляемость: контрольные точки в отдельной коллекции позволяют
    остановить робота в любой момент и продолжить с того же места при
    перезапуске;
    \item переобкачка по изменению: сравнение MD5-хэша содержимого позволяет
    не трогать неизменившиеся документы и сохранять историю версий изменившихся;
    \item реальность снова внесла поправки: источники оказались JSON API, а не
    HTML-сайтами, а robots.txt MusicBrainz формально запрещает доступ к
    собственному публичному API --- пришлось делать проверку robots.txt
    настраиваемой для каждого источника.
\end{itemize}

Скорость обкачки упирается не в сеть, а в задержку между запросами (2~с):
полный проход по корпусу из 30000 документов занимает почти сутки. Для
следующих лабораторных работ в базе теперь есть актуальные документы с датами
обкачки и историей версий, на которых можно строить поисковый индекс.

\pagebreak
